Image recognition is a fundamental topic in artificial intelligence and computer vision, with broad applications in intelligent security, assisted medical diagnosis, industrial inspection, and autonomous driving. Benefiting from the rapid development of deep learning, convolutional neural networks have achieved remarkable performance improvements over handcrafted feature based methods. Nevertheless, existing models still suffer from large parameter sizes, high deployment cost, and insufficient robustness in complex scenarios. Therefore, it is meaningful to investigate a balanced solution that improves recognition accuracy while preserving computational efficiency.

This thesis focuses on the design of a lightweight image recognition model. First, the principles of convolutional neural networks, residual learning, and attention mechanisms are reviewed, and the strengths and weaknesses of representative image recognition methods are summarized. Second, an improved architecture that combines channel re-calibration and multi-scale feature aggregation is constructed on top of a baseline convolutional network. The proposed design enhances feature representation by encouraging information interaction between shallow and deep layers. Third, a complete experimental workflow is developed by integrating data augmentation, transfer learning, and ablation studies. The model is evaluated from the perspectives of classification accuracy, parameter size, and inference latency.

Experimental results indicate that the proposed method consistently outperforms the baseline network on public datasets while maintaining a relatively small model size and low inference overhead. Visualization of misclassified samples further shows that multi-scale feature fusion helps alleviate the negative influence of cluttered backgrounds, scale variation, and partial occlusion. The study demonstrates that a well organized undergraduate thesis can effectively present problem formulation, method design, experimental verification, and academic writing in a coherent manner.

\gxuenkeywords{Image recognition; Deep learning; Convolutional neural network; Lightweight model; Attention mechanism}
